\documentclass[a4paper, serif, 10pt]{moderncv}

%% moderncv
% color and style
\moderncvcolor{black}
\moderncvstyle{banking}

%% page layout/margins
\usepackage[top=2.5cm, bottom=2.5cm, left=1.5cm, right=1.5cm]{geometry}

\nopagenumbers{}

%% Babel config for Indonesian (Bahasa Indonesia) language
% ref:
% - https://latex3.github.io/babel/guides/locale-indonesian.html
\usepackage[indonesian]{babel}

%% fontspec
\usepackage{fontspec}

\IfFontExistsTF{Linux Libertine}{
  \setmainfont[Ligatures={TeX, Common}, Numbers=OldStyle]{Linux Libertine}
}{}
\IfFontExistsTF{Linux Libertine O}{
  \setmainfont[Ligatures={TeX, Common}, Numbers=OldStyle]{Linux Libertine O}
}{}

%% CTeX
% CJK support, used to show CJK characters in the resume
%
% - heading=false: disable ctex heading
% - punct=kaiming: use kaiming punctuations styles for CJK
% - scheme=plain: use plain scheme, do not override \normalsize font size
% - space=auto: space settings for CJK characters
%
% CJK fonts are resolved per locale language, e.g. Noto Serif/Sans CJK JP
% for Japanese resumes, with fallback to the next region if not installed
%
% ref:
% - http://ctan.mirrorcatalogs.com/language/chinese/ctex/ctex.pdf
\usepackage[UTF8, heading=false, punct=kaiming, scheme=plain, space=auto]{ctex}

\IfFontExistsTF{Noto Serif CJK SC}{
  \setCJKmainfont{Noto Serif CJK SC}
}{}
\IfFontExistsTF{Noto Sans CJK SC}{
  \setCJKsansfont{Noto Sans CJK SC}
}{}

%% line spacing
\usepackage{setspace}
\setstretch{1.125}

%% URL styling - use normal text instead of monospace
\urlstyle{same}

%% Patch \httplink and \httpslink to handle full URLs
% The original moderncv macros always prepend http:// or https:// to URLs.
% This causes issues when the URL already has a protocol (e.g., https://example.com)
% resulting in malformed URLs like https://https://example.com.
% This patch checks if the URL already contains "://" and uses it directly if so.
% Note: We use \str_set:Nx (x-expansion) to expand macros like \@homepage first.
\ExplSyntaxOn
\str_new:N \g_yamlresume_url_str

\RenewDocumentCommand{\httpslink}{O{}m}{
  \str_set:Nx \g_yamlresume_url_str {#2}
  \str_if_in:NnTF \g_yamlresume_url_str {://}
    { \url{#2} }
    { \url{https://#2} }
}

\RenewDocumentCommand{\httplink}{O{}m}{
  \str_set:Nx \g_yamlresume_url_str {#2}
  \str_if_in:NnTF \g_yamlresume_url_str {://}
    { \url{#2} }
    { \url{http://#2} }
}
\ExplSyntaxOff

\name{Anindya Pratiwi}{}
\title{Perawat Terdaftar dengan Pengalaman Rawat Inap dan Kegawatdaruratan}
\phone[mobile]{+62 812 3456 7890}
\email{anindya.pratiwi@email.com}
\homepage{https://www.anindyapratiwi.id}

\address{Jl. Dharmahusada Indah No. 45, Surabaya, Jawa Timur, Indonesia, 60115}{}{}

\extrainfo{{\small \faLinkedin}\ \href{https://www.linkedin.com/in/anindya-pratiwi}{@anindya-pratiwi} {} {} {} • {} {} {} 
{\small \faInstagram}\ \href{https://www.instagram.com/anindya.pratiwi}{@anindya.pratiwi}}

\begin{document}

\maketitle

\section{Informasi Dasar}

\cvline{}{\begin{itemize}
\item Perawat terdaftar dengan pengalaman lebih dari enam tahun di rumah sakit dan klinik.
\item Terampil dalam asuhan keperawatan rawat inap, triage, manajemen obat, dan dokumentasi rekam medis elektronik.
\item Berkomitmen pada keselamatan pasien, pencegahan infeksi, dan komunikasi terapeutik.
\item Mampu bekerja dalam tim multidisiplin serta membimbing perawat baru dan mahasiswa keperawatan.
\end{itemize}}
\section{Pendidikan}

\cventry{Sep 2013–Agu 2017}
        {Sarjana, Ilmu Keperawatan, Nilai: 3.72}
        {Universitas Kesehatan Nusantara}
        {\url{https://www.ukn.ac.id}}
        {}
        {\begin{itemize}
\item Menyelesaikan pendidikan sarjana keperawatan dengan predikat cum laude.
\item Menjalani praktik klinik di rumah sakit pendidikan dan puskesmas.
\item Aktif dalam organisasi kemahasiswaan bidang kesehatan masyarakat.
\end{itemize}
\textbf{Mata Kuliah}: Asuhan Keperawatan Dewasa, Keperawatan Gawat Darurat, Farmakologi Klinik, Keperawatan Anak, Keperawatan Maternitas, Manajemen Keperawatan, Promosi Kesehatan, Biostatistika}
\section{Pengalaman Kerja}

\cventry{Jan 2021–Sekarang}
        {Perawat Terdaftar Rawat Inap}
        {RS Prima Sehat Surabaya}
        {\url{https://www.rsprimasehat.co.id}}
        {}
        {\begin{itemize}
\item Memberikan asuhan keperawatan komprehensif kepada 20-25 pasien rawat inap per shift.
\item Melakukan pengkajian awal, pemantauan tanda vital, dan dokumentasi pada sistem rekam medis elektronik.
\item Mengelola pemberian obat, cairan infus, dan tindakan keperawatan sesuai standar prosedur.
\item Berkolaborasi dengan dokter, apoteker, dan tim multidisiplin dalam perencanaan pulang pasien.
\item Menjadi preceptor bagi perawat baru dan mahasiswa keperawatan.
\end{itemize}
\textbf{Kata Kunci}: Rawat inap, Asuhan keperawatan, Rekam medis elektronik, Keselamatan pasien, Pembimbing klinik}

\cventry{Sep 2017–Des 2020}
        {Perawat Pelaksana}
        {Klinik Utama Cahaya Medika}
        {\url{https://www.klinikcahayamedika.co.id}}
        {}
        {\begin{itemize}
\item Memberikan pelayanan keperawatan rawat jalan dan tindakan minor.
\item Melakukan triage awal, pengambilan sampel darah, dan pemeriksaan penunjang dasar.
\item Mengedukasi pasien dan keluarga tentang pengelolaan penyakit kronis dan obat.
\item Menjaga kebersihan dan sterilisasi alat sesuai prinsip pencegahan infeksi.
\item Mendukung program vaksinasi dan pemeriksaan kesehatan berkala.
\end{itemize}
\textbf{Kata Kunci}: Rawat jalan, Triage, Edukasi pasien, Pencegahan infeksi, Vaksinasi}
\section{Bahasa}

\cvline{Indonesia}{Bahasa Ibu / Bilingual \hfill \textbf{Kata Kunci}: Bahasa ibu}
\cvline{Inggris}{Tingkat Profesional Penuh \hfill \textbf{Kata Kunci}: TOEFL 580, Komunikasi klinis}
\section{Keahlian}

\cvline{Keperawatan Klinis}{Ahli \hfill \textbf{Kata Kunci}: Asuhan keperawatan, Pemantauan tanda vital, Manajemen obat, Perawatan luka, Infus dan kateter}
\cvline{Kegawatdaruratan}{Mahir \hfill \textbf{Kata Kunci}: BLS, ACLS, Triage, Resusitasi}
\cvline{Manajemen Pasien}{Mahir \hfill \textbf{Kata Kunci}: Perencanaan pulang, Koordinasi tim, Edukasi pasien}
\cvline{Sistem Informasi Kesehatan}{Menengah \hfill \textbf{Kata Kunci}: Rekam medis elektronik, SIMRS, Dokumentasi}
\cvline{Keselamatan Pasien}{Mahir \hfill \textbf{Kata Kunci}: Pencegahan infeksi, Hand hygiene, Identifikasi pasien, Pelaporan insiden}
\section{Penghargaan}

\cventry{Des 2023}
        {RS Prima Sehat Surabaya}
        {Perawat Teladan Tahun 2023}
        {}
        {}
        {\begin{itemize}
\item Diberikan atas dedikasi tinggi dalam asuhan keperawatan dan kolaborasi tim.
\item Diakui karena inisiatif peningkatan kepatuhan hand hygiene di ruang rawat inap.
\end{itemize}}

\cventry{Agu 2017}
        {Universitas Kesehatan Nusantara}
        {Cum Laude}
        {}
        {}
        {\begin{itemize}
\item Meraih predikat cum laude dengan IPK 3,72.
\end{itemize}}
\section{Sertifikat}

\cventry{Feb 2018}
        {Kementerian Kesehatan Republik Indonesia}
        {Surat Tanda Registrasi (STR) Perawat}
        {\url{https://www.kemkes.go.id}}
        {}
        {}

\cventry{Mei 2023}
        {American Heart Association}
        {Sertifikat Bantuan Hidup Dasar (BLS)}
        {\url{https://www.heart.org}}
        {}
        {}

\cventry{Jun 2023}
        {American Heart Association}
        {Sertifikat Advanced Cardiovascular Life Support (ACLS)}
        {\url{https://www.heart.org}}
        {}
        {}

\cventry{Sep 2022}
        {RS Prima Sehat Surabaya}
        {Pelatihan Pencegahan dan Pengendalian Infeksi}
        {}
        {}
        {}
\section{Publikasi}

\cventry{Nov 2022}
        {Hubungan Beban Kerja Perawat dengan Kualitas Asuhan Keperawatan di Ruang Rawat Inap}
        {Jurnal Keperawatan Nusantara}
        {\url{https://jurnal.keperawatannusantara.id/index.php/jkn/article/view/123}}
        {}
        {\begin{itemize}
\item Studi cross-sectional terhadap 84 perawat di dua rumah sakit swasta di Surabaya.
\item Menemukan korelasi signifikan antara beban kerja dan kepatuhan dokumentasi asuhan.
\end{itemize}}
\section{Proyek}

\cventry{Feb 2023–Agu 2023}
        {Proyek perbaikan mutu keperawatan berbasis data untuk menurunkan risiko infeksi nosokomial.}
        {Peningkatan Kepatuhan Hand Hygiene}
        {\url{https://www.rsprimasehat.co.id/mutu-hand-hygiene}}
        {}
        {\begin{itemize}
\item Memimpin proyek perbaikan mutu untuk meningkatkan kepatuhan hand hygiene di ruang rawat inap.
\item Melakukan audit kepatuhan, pelatihan singkat, dan penyediaan pengingat visual.
\item Meningkatkan kepatuhan dari 68\% menjadi 89\% dalam enam bulan.
\end{itemize}
\textbf{Kata Kunci}: Keselamatan pasien, Hand hygiene, Perbaikan mutu, Pencegahan infeksi}
\section{Minat}

\cvline{Kesehatan Masyarakat}{Posyandu, Edukasi gizi, Donor darah}
\cvline{Pengembangan Diri}{Membaca jurnal, Pelatihan klinis, Musik}
\section{Kegiatan Sukarela}

\cventry{Jan 2019–Des 2023}
        {Perawat Sukarelawan}
        {Posyandu Melati Surabaya}
        {\url{https://www.posyandu-melati.id}}
        {}
        {\begin{itemize}
\item Membantu pemeriksaan kesehatan ibu dan anak serta imunisasi dasar.
\item Memberikan edukasi gizi dan pencegahan stunting kepada keluarga.
\item Mendukung pencatatan tumbuh kembang anak secara berkala.
\end{itemize}}
    

\cventry{Mar 2018–Des 2021}
        {Relawan Kesehatan}
        {Palang Merah Indonesia Kota Surabaya}
        {\url{https://www.pmi.or.id}}
        {}
        {\begin{itemize}
\item Terlibat dalam kegiatan donor darah dan pemeriksaan kesehatan gratis.
\item Memberikan pertolongan pertama pada kegiatan komunitas.
\end{itemize}}
    
\end{document}